\section{Design Principles}\label{sec:6}

Design principles are guidelines for developers. Unlike architecture
decisions, they guide judgement rather than impose a rule.

\subsection{SOLID Principles}\label{sec:6.1}

{\sloppy The table names the classes in the repository that apply each principle
(modules \texttt{herd-health}, \texttt{access-policy} and
\texttt{farm-edge} under \texttt{backend/}). \Cref{sec:m2-health-solid} explains the
Health authorization slice in more detail. The M1 report used
design-level names (\texttt{MovementRecorder}, \texttt{HeadCountCalculator},
\texttt{AlertRule}, \texttt{ReportExporter}, \texttt{ReaderEventSource},
\texttt{ReaderHealthMonitor}); the M2 prototype does not implement them, so
they are replaced here by the classes that exist.\par}

{\small
\begin{longtable}[]{|>{\raggedright\arraybackslash}p{\dimexpr 0.2280\linewidth-2\tabcolsep\relax}|>{\raggedright\arraybackslash}p{\dimexpr 0.5092\linewidth-2\tabcolsep\relax}|>{\raggedright\arraybackslash}p{\dimexpr 0.2328\linewidth-2\tabcolsep\relax}|}
\hline
\textbf{Principle} & \textbf{Application in VinaBeef (code)} & \textbf{Supports} \\ \hline
\endhead
Single Responsibility & \texttt{HealthController} handles HTTP only;
\texttt{HealthIntentService} runs the use cases; \texttt{HealthAuthorization}
applies the authorization decision order; \texttt{JdbcHealthRepository}
owns the SQL. At the farm edge, \texttt{PassageClassifier} is a pure
classification rule (no Spring, no database) and \texttt{GateService} owns
the transaction. & NFR-08 \\ \hline
Open/Closed & A new identity or grant source is a new bean behind
\texttt{TrustedPrincipalProvider} or \texttt{CurrentGrantContextProvider};
\texttt{HealthBoundaryConfiguration} registers fail-closed defaults with
\texttt{@ConditionalOnMissingBean}, so the use cases do not change. &
NFR-04, NFR-08 \\ \hline
Liskov Substitution & Every \texttt{CurrentGrantContextProvider} (the JSON
fixture \texttt{JsonFixtureCurrentGrantContextProvider} and the fail-closed
default) either returns the grant context for the verified subject or fails
as unavailable; an unavailable source never yields an allow. & NFR-04 \\ \hline
Interface Segregation & Five inbound ports: \texttt{LookupCattleUseCase},
\texttt{SubmitHealthIntentUseCase}, \texttt{GetHealthIntentReceiptUseCase},
\texttt{GetHealthHistoryUseCase}, \texttt{ListCattleInScopeUseCase}; a
client such as \texttt{HealthController} depends only on the operations it
calls. & NFR-08 \\ \hline
Dependency Inversion & \texttt{HealthIntentService} depends on
\texttt{HealthRepository}, implemented by \texttt{JdbcHealthRepository};
\texttt{HealthAuthorization} depends on \texttt{AccessPolicyPort},
implemented by \texttt{ClinicalAuthorizationPolicyAdapter}, which calls
\texttt{ClinicalAuthorizationPolicy} in \texttt{libs/access-policy} (plain
Java, no Spring). & NFR-04, NFR-08 \\ \hline
\end{longtable}
}

\subsection{Other Principles}\label{sec:6.2}


\begin{itemize}
  \item Complete at the edge, consolidate at the centre: nothing on the gate path waits on the WAN (ADR-02, ADR-11).
  \item Persist, then acknowledge: never report success for work that is not durably committed (ADR-03).
  \item Assume duplicates: any message or request may arrive more than once; never assume a remote receiver is idempotent unless its contract says so (ADR-04, ADR-06).
  \item Deny by default, least privilege: each role, service and machine identity gets only the permissions it needs (ADR-07).
  \item Evidence with the change: an audited operation and its evidence commit together (ADR-08).
  \item Loose coupling, high cohesion: communicate through events and versioned APIs; keep logic that changes together in one service (ADR-09).
  \item Observability: structured logs, metrics and reader heartbeats; every log and audit record carries an operation or correlation ID.
  \item Keep it simple: add a service or a technology only when a driving characteristic requires it.
\end{itemize}

\subsection{Applied SOLID in Health Authorization}\label{sec:m2-health-solid}

\textbf{Single Responsibility.} The Java policy evaluates grants against
the action and resource supplied by Health. The application service
enforces the decision and coordinates acceptance, retry and audit
evidence within the Health transaction. A change to a scope rule can
be exercised in policy tests. A change to storage or transaction
behavior is exercised through the Health backend.

\textbf{Dependency Inversion.} Health's application layer calls
\texttt{AccessPolicyPort}. The concrete adapter maps Health's principal,
grant and resource types to the library. The policy rules are testable
without Spring or PostgreSQL. Running the integrated slice verifies
that the adapter supplies the right context and Health enforces the
decision. \Cref{sec:m2-health-auth,sec:m2-health-results} describes both
levels of verification.

\subsection{Reflection on the Health Slice}\label{sec:m2-health-reflection}

The PostgreSQL suite exposed a failure that the policy unit tests could
not exercise. The JDBC driver could not bind a Java \texttt{Instant}
directly to a timestamp parameter during a clinical write. Converting
the timestamp to an offset date-time at the JDBC boundary preserved
the domain model and the authorization contract. The rollback test
also checks that its exception came from the injected source-evidence
insert failure before it checks the database state. That assertion
matters because an unrelated error early in the write could also
leave zero rows. These two checks show why policy unit tests and
database integration tests both belong in the verification plan.
